\documentclass[11pt]{article}
\usepackage[a4paper,margin=20mm]{geometry}
\usepackage{amsmath,amssymb,graphicx}
\graphicspath{{../../}}
\setlength{\parindent}{0pt}
\setlength{\emergencystretch}{2em}
\newcommand{\biaspage}[4]{%
  \begin{center}\Large #1: #2\end{center}
  \includegraphics[width=\linewidth]{figs/record_length_aggregate_bias_#3.pdf}
  \par\medskip
  Each panel shows the estimated aggregate bias
  \[
    \widehat B_N=\|\widehat{\boldsymbol b}_N\|_2
      =\left(\sum_{j=1}^{6}\widehat b_{j,N}^{\,2}\right)^{1/2},
    \qquad
    \widehat b_{j,N}=\frac{1}{R_N}\sum_{r\in\mathcal I_N}
      (\widehat\theta_{j,N}^{(r)}-\theta_j^\star).
  \]
  All six parameters enter on their original scale, with equal weights.
  This is the norm of the mean-error vector; the mean norm of the
  individual estimation errors is a different quantity.
  Finite stalled and iteration-limited outputs are retained; failed fits
  are excluded. No Monte Carlo intervals are plotted for this aggregate.
  Linear vertical axes start at zero and are capped at one, with tighter
  ranges retained. Boundary triangles mark values above the display cap.
  #4
  \clearpage}
\begin{document}
\biaspage{Example 1}{PL}{example_1_pl_recursive_PL}{}
\biaspage{Example 1}{FSM, 128 coefficients}{example_1_fsm_fsm_128}%
  {For GHAM-2-I, the retained counts are shown in record-length order;
   the final point summarizes only 7 of 100 runs.}
\biaspage{CSTR}{PL}{example_CSTR_pl_recursive_PL}{}
\biaspage{CSTR}{FSM, 128 coefficients}{example_CSTR_fsm_fsm_128}{}
\biaspage{Example 1}{FSM, 256 coefficients}{example_1_fsm_fsm_256}%
  {For GHAM-2-I, the final point summarizes only 7 of 100 runs.}
\biaspage{CSTR}{FSM, 256 coefficients}{example_CSTR_fsm_fsm_256}{}
\end{document}
